%+------------------------------------------------------------------------+
%| NOTE: LaTeX source: 2026-Cv-Tony-DeepBlue-DataAnalystModeller.
%| DATE: 2026-11-04
%| Author: Antonio Michele Miti
%| Credits: Original source and provenance notes retained below.
%+------------------------------------------------------------------------+
% CV approvato il 04/11/2026; PDF definitivo senza colori di revisione.
% Preparato con assistenza AI; provenienza conservata nelle note della candidatura.
% Preparato con assistenza AI; provenienza conservata nelle note della candidatura.
% Competenze documentate, bibliografia selettiva, verifica estrazione testo.
\documentclass[10pt,a4paper]{article}


%=======================================================================%
%      1. PREAMBOLO - PACCHETTI E IMPAGINAZIONE
%=======================================================================%

\usepackage[utf8]{inputenc}
\usepackage[T1]{fontenc}
\usepackage[italian]{babel}
\usepackage{lmodern}
\usepackage[margin=1.7cm,headheight=14pt]{geometry}
\usepackage{xcolor}
\usepackage{enumitem}
\usepackage{tabularx}
\usepackage{microtype}
\usepackage{needspace}
\usepackage{fancyhdr}
\usepackage{lastpage}
\usepackage{fontawesome5}

\usepackage{hyperref}

\definecolor{accent}{HTML}{000000}
\definecolor{muted}{HTML}{000000}

\hypersetup{
    colorlinks=true,
    urlcolor=accent,
    linkcolor=accent,
    pdfauthor={Antonio Michele Miti},
    pdftitle={Antonio Michele Miti - Deep Blue Data analyst and modeller}
}

\setlength{\parindent}{0pt}
\setlength{\parskip}{0pt}
\setlength{\emergencystretch}{2em}
\setlist[itemize]{
    leftmargin=1.3em,
    itemsep=2pt,
    topsep=3pt,
    parsep=0pt,
    partopsep=0pt
}

\pagestyle{fancy}
\fancyhf{}
\renewcommand{\headrulewidth}{0pt}
\cfoot{%
    \footnotesize\textcolor{muted}{%
        Antonio Michele Miti \enspace | \enspace
        Pagina \thepage\ di \pageref{LastPage}%
    }%
}


%=======================================================================%
%      2. MACRO - ELEMENTI RIUTILIZZABILI
%=======================================================================%

% Titolo di sezione con filetto e spazio minimo per il contenuto.
% Usage: \cvsection{Title}
\newcommand{\cvsection}[1]{%
    \Needspace{4\baselineskip}%
    \vspace{0.8em}%
    {\large\bfseries\color{accent} #1}\par
    \vspace{0.2em}\hrule\vspace{0.45em}%
}

% Voce con titolo, data, istituzione e descrizione.
% Usage: \entry{Title}{Date}{Institution}{Description}
\newcommand{\entry}[4]{%
    \Needspace{4\baselineskip}%
    \textbf{#1}\hfill #2\par
    \textit{#3}\par
    #4\par\vspace{0.45em}%
}

% Titolo di studio con voto sotto la data, nella colonna destra.
% Usage: \education{Title}{Date}{Grade}{Institution}{Description}
\newcommand{\education}[5]{%
    \Needspace{5\baselineskip}%
    \begin{tabularx}{\textwidth}{@{}>{\raggedright\arraybackslash}Xr@{}}
        \textbf{#1} & #2 \\
        \textit{#4} & #3
    \end{tabularx}\par
    #5\par\vspace{0.45em}%
}

% Progetto con rimando alle evidenze e lista dei contributi.
% Usage: \project{Title}{Evidence link}{Items}
\newcommand{\project}[3]{%
    \Needspace{5\baselineskip}%
    \textbf{#1}\par
    #2\par
    \begin{itemize}
        #3
    \end{itemize}
    \vspace{0.25em}%
}


%=======================================================================%
%      3. DOCUMENTO - PROFILO E PERCORSO PROFESSIONALE
%=======================================================================%

% Unicode mapping and explicit interword spaces for PDF text extraction.
\input{glyphtounicode}
% Text equivalents for decorative contact icons.
\pdfglyphtounicode{envelope}{0045 006D 0061 0069 006C 003A}
\pdfglyphtounicode{phone}{0054 0065 006C 003A}
\pdfglyphtounicode{globe}{0057 0065 0062 003A}
\pdfglyphtounicode{github}{0047 0069 0074 0048 0075 0062 003A}
\pdfglyphtounicode{orcid}{0020}
\pdfgentounicode=1
\pdfinterwordspaceon
\begin{document}
{\Huge\bfseries Antonio Michele Miti}\par
\vspace{0.25em}
{\large\color{accent} Fisica matematica e calcolo scientifico}\par
    \vspace{0.1em}
    \begin{center}
        \href{mailto:antoniomichele.miti@gmail.com}{%
                \faEnvelope\enspace antoniomichele.miti@gmail.com}  |
        \href{tel:+393339825418}{\faPhone\enspace +39 333 9825 418} \\[3pt]
        \href{https://www.antoniomiti.it/}{\faGlobe\enspace antoniomiti.it} |
        \href{https://github.com/MasterToninus}{\faGithub\enspace MasterToninus} |
        \href{https://orcid.org/0000-0002-8829-1943}{%
            \faOrcid\enspace ORCID: 0000-0002-8829-1943}
    \end{center}
    \vspace{0.1em}
\textit{Candidatura: Data analyst and modeller}\par
\textit{Deep Blue -- Roma / Ibrido}\par



\cvsection{Profilo}
{Ricercatore in fisica matematica, con laurea magistrale in Fisica e dottorato internazionale in Matematica. Intendo trasferire le competenze maturate nella ricerca alla modellazione scientifica in ambito industriale, lavorando su problemi di domini diversi. La preparazione matematica e fisica, insieme all'esperienza di ricerca e ai progetti universitari di calcolo scientifico, costituisce la base per acquisire rapidamente i metodi e gli strumenti delle nuove applicazioni.}

\cvsection{Competenze pertinenti}
\begin{itemize}
    \item {\textbf{Metodi matematici:} geometria differenziale e fisica matematica; osservabili, simmetrie, vincoli e riduzione di sistemi teorici.}
    \item {\textbf{Calcolo scientifico:} implementazione e confronto di algoritmi, test, prestazioni e stabilità numerica in progetti universitari C++/CUDA.}
    \item {\textbf{Programmazione e strumenti:} Python nel percorso di formazione scientifica; Java in uno stage formativo del 2016; Mathematica, Git e \LaTeX.}
    \item {\textbf{Ricerca e comunicazione:} autonomia nell'impostazione dei problemi, documentazione scientifica, pubblicazioni e didattica in contesti internazionali; uso professionale dell'inglese.}
\end{itemize}

\cvsection{Esperienza professionale nella ricerca}
\entry{Postdoctoral Fellow -- CIVIS3i, MSCA COFUND}{2024--2026}{Sapienza Università di Roma}{%
Ricerca nel progetto MultiSymORI su geometria multisimplettica, osservabili e riduzione, nell'ambito della fisica matematica.
Curato il lavoro scientifico, le collaborazioni internazionali, le visite di ricerca, le pubblicazioni e la rendicontazione del progetto.
}
\entry{Assegnista di ricerca}{2022--2024}{Università Cattolica del Sacro Cuore, Brescia}{%
Ricerca in geometria simplettica e multisimplettica con applicazioni alla fisica matematica; collaborazione internazionale e didattica universitaria.}
\entry{Research Fellow}{2021--2022}{Max Planck Institute for Mathematics, Bonn}{%
Ricerca sulle strutture algebriche superiori in geometria e fisica matematica, in un ambiente internazionale.}

\cvsection{Formazione}
\education{Dottorato in Science: Mathematics}{2021}{}{Università Cattolica del Sacro Cuore / KU Leuven}{%
Dottorato internazionale in Matematica (2016--2021), con periodi di ricerca in Italia e Belgio; specializzazione in geometria differenziale e fisica matematica.}
\education{Laurea magistrale in Fisica -- curriculum teorico}{2015}{110/110}{Università degli Studi di Milano}{%
Metodi matematici della fisica, meccanica geometrica e metodi computazionali.}
\education{Laurea in Fisica}{2010}{99/110}{Università degli Studi di Milano-Bicocca}{%
Tesi sulle applicazioni dei gruppi di Lie alla meccanica del corpo rigido.}

\cvsection{Lingue}
\begin{itemize}
    \item \textbf{Italiano:} madrelingua.
    \item \textbf{Inglese:} uso professionale nella ricerca e nella didattica; TOEFL iBT 102/120 (2018, risultato storico).
    \item \textbf{Francese:} base.
\end{itemize}

\newpage
\cvsection{Progetti di calcolo scientifico}
\entry{Simulazione Monte Carlo del modello di Ising}{Progetto universitario}{C++}{%
\begin{itemize}
    \item Implementato e confrontato gli algoritmi di Metropolis e a cluster mediante esperimenti numerici sullo stesso sistema fisico.
\end{itemize}
\href{https://github.com/MasterToninus/Computational-Physics-Exercise/tree/master/Ising\%20simulation\%20Metropolis\%20vs\%20Cluster}{Codice del progetto}.}
\entry{Calcolo di determinanti su CPU e GPU}{Progetto universitario}{C++ / CUDA}{%
\begin{itemize}
    \item Implementato un algoritmo di condensazione con una libreria di matrici CUDA, kernel GPU, test e strategie di pivoting.
    \item Confrontato l'implementazione GPU con quella sequenziale CPU al variare delle dimensioni delle matrici, studiando accessi alla memoria, prestazioni e stabilità numerica.
\end{itemize}
\href{https://github.com/MasterToninus/Computational-Physics-Exercise/tree/master/Condensation\%20Algorithm\%20in\%20CUDA}{Codice del progetto}.}

\cvsection{Esperienze di supporto informatico}
\entry{Supporto ai sistemi informatici dipartimentali}{2017--2018}{Università Cattolica del Sacro Cuore, Brescia}{%
Attività di assistenza al mantenimento dei sistemi informatici del Dipartimento di Matematica e Fisica, durante il dottorato.}
\entry{Stage di programmazione Java}{2016}{CST s.r.l., Milano}{%
Formazione in Java, programmazione orientata agli oggetti e strumenti per lo sviluppo software professionale.}
\entry{Collaborazione studentesca -- mantenimento cluster}{2014--2015}{Laboratorio di Calcolo e Multimedia, Università degli Studi di Milano}{%
Mantenimento del cluster del laboratorio e supporto agli utenti e alle attività didattiche di fisica computazionale in ambiente Linux.}

\cvsection{Aggiornamento professionale}
\begin{itemize}
    \item \textbf{2026} -- Iscritto alla selezione del Master universitario di II livello in \emph{High-Performance and Quantum Computing}, Università di Bologna, a.a.\ 2026/2027 (codice 6162).
    \item \textbf{2024} -- Workshop \emph{High Performance Computing and Quantum Computing}, CINECA.
    \item \textbf{2018} -- \emph{Scientific Computing with Python}, Università Cattolica.
    \item \textbf{2017} -- \emph{Python for Computational Science}, CINECA-SCAI.
    \item \textbf{2016} -- \emph{Introduction to Scientific and Technical Computing in C++}, CINECA-SCAI.
\end{itemize}

\cvsection{Pubblicazioni selezionate}
\begin{itemize}[itemsep=4pt]
    \item A. Luongo, A. M. Miti, V. Narasimhachar, A. Sireesh.
    \emph{Measurement-based uncomputation of quantum circuits for modular arithmetic}.
    DAC (2025).
    \href{https://doi.org/10.1109/DAC63849.2025.11132981}{DOI}.
    \item A. M. Miti, L. Ryvkin.
    \emph{Multisymplectic observable reduction using constraint triples}.
    Differential Geometry and its Applications \textbf{100} (2025), 102272.
    \href{https://doi.org/10.1016/j.difgeo.2025.102272}{DOI}.
    \item C. Blacker, A. M. Miti, L. Ryvkin.
    \emph{Reduction of $L_\infty$-algebras of observables on multisymplectic manifolds}.
    SIGMA \textbf{20} (2024), 061.
    \href{https://doi.org/10.3842/SIGMA.2024.061}{DOI}.
    \item A. M. Miti, M. Spera.
    \emph{A hydrodynamical homotopy co-momentum map and a multisymplectic interpretation of higher-order linking numbers}.
    Journal of the Australian Mathematical Society (2021, pubblicazione online).
    \href{https://doi.org/10.1017/S1446788720000518}{DOI}.
\end{itemize}
\href{https://orcid.org/0000-0002-8829-1943}{Profilo scientifico e pubblicazioni: ORCID}.\par

\cvsection{Attività scientifica e didattica in numeri}
\begin{minipage}[t]{0.48\textwidth}
\textbf{Ricerca}\par\vspace{3pt}
\begin{tabularx}{\linewidth}{|>{\raggedright\arraybackslash}X|r|}
\hline
Pubblicazioni scientifiche & 7 \\ \hline
Preprint non ancora pubblicati & 1 \\ \hline
Lavori in preparazione & 3 \\ \hline
Anni di esperienza postdoc (circa) & 5 \\ \hline
\end{tabularx}
\end{minipage}\hfill
\begin{minipage}[t]{0.48\textwidth}
\textbf{Comunicazione e attività internazionale}\par\vspace{3pt}
\begin{tabularx}{\linewidth}{|>{\raggedright\arraybackslash}X|r|}
\hline
Interventi su invito & 17 \\ \hline
Comunicazioni e poster & 20 \\ \hline
Conferenze frequentate & 62 \\ \hline
Soggiorni di ricerca & 12 \\ \hline
\end{tabularx}
\end{minipage}
\par\vspace{5pt}
\textbf{Didattica e organizzazione:} 5 corsi di livello magistrale/dottorale, 58 ore di didattica; 4 conferenze organizzate.

\vspace{0.5em}
{\footnotesize{Autorizzo Deep Blue S.r.l. al trattamento dei dati personali contenuti nel presente curriculum per le finalità connesse alla selezione del personale, ai sensi del Regolamento UE 2016/679.}}\par

\end{document}
